\subsection{標準}
標準は、比較のための基準、許容範囲を示す分類、要求される水準などを定めるものである。工学では、製品、過程、材料が受け入れ可能な品質を持つように、要求、仕様、指針を提供する。
標準への適合は、組織の製品や過程が一定の要求を満たすことを社会に示す手段になる。ただし、標準が有用であるためには、適合することが実際または知覚上の価値を生む必要がある。
標準は、購入者を保護し、事業者を保護し、ソフトウェア工学で使う方法や手順を明確にし、共通用語と期待を提供する。国際的な標準化組織には、ITU、IEC、IEEE、ISOなどがある。国や地域にも標準化組織がある。
標準は、多くの場合、合意形成に基づいて作られる。技術者は、自分の領域で適用される標準を把握し、標準の変更にも追随する必要がある。法律や契約が特定標準への適合を求める場合、標準は設計制約の最小要求になる。
\par\smallskip
\begin{center}
\centering
\IfFileExists{assets/generated_figures/ch18/fig-ch18-13.png}{\includegraphics[width=0.96\linewidth]{assets/generated_figures/ch18/fig-ch18-13.png}}{\fbox{\parbox[c][48mm][c]{0.9\linewidth}{\centering 画像生成待ち\\標準と適合の概念図}}}
\par\smallskip
\refstepcounter{figure}\label{fig:ch18-13}図 \thefigure: 標準と適合の概念図
\end{center}
図 \thefigure は、標準と適合の概念図を視覚的に整理したものである。
\par\smallskip
